% Requires: \usetikzlibrary{positioning,arrows.meta}
\begin{figure}[H]
  \centering
  \resizebox{0.97\textwidth}{!}{%
  \begin{tikzpicture}[
      >={Stealth[length=1.8mm]},
      res/.style={draw=black!70,line width=0.4pt,fill=black!6,
        rounded corners=2pt,align=center,inner sep=3.5pt,
        font=\footnotesize},
      flow/.style={draw=black,line width=0.5pt,->}
    ]
    \node[res,text width=2.1cm] (D) at (0,0)
      {design $(y,z)$\\
       \textcolor{black!60}{\scriptsize from master or enumeration}};
    \node[res,text width=2.6cm] (O) at (3.6,0)
      {routing oracle\\
       \textcolor{black!60}{\scriptsize one transportation block per
       product $l$}};
    \node[res,text width=3.0cm] (C) at (7.5,0)
      {primal and dual certificates\\
       \textcolor{black!60}{\scriptsize optimal flow and
       $u^\ast\in\Delta_l$, or a ray}};
    \node[res,text width=2.6cm] (K) at (11.2,0)
      {disaggregated cut\\
       \textcolor{black!60}{\scriptsize optimality on $\theta_l$, or
       feasibility}};
    \draw[flow] (D) -- (O);
    \draw[flow] (O) -- (C);
    \draw[flow] (C) -- (K);
  \end{tikzpicture}%
  }
  \caption{From a design to a cut, the pipeline of
  Proposition~\ref{prop:duals}. Each oracle call of
  Proposition~\ref{prop:oracle} prices the design one product at a
  time, returns an optimal flow together with a dual point
  $u^\ast\in\Delta_l$ certifying the block value, and that dual point
  yields the disaggregated Benders optimality cut, valid for every
  design with a feasible block because $\Delta_l$ never depends on $(y,z)$, while an
  infeasible block yields a recession ray and a feasibility cut
  instead.}
  \label{fig:bendersflow}
\end{figure}
